\chapter{{\fontsize{14pt}{21pt}\selectfont\MakeUppercase{Тестирование и оценка результатов}}}
\label{cha:testing}

\section{Цель и программа испытаний}

Цель тестирования состоит в проверке того, что веб-приложение выполняет заявленные функции и не выходит за пределы программной реализации. Проверка ориентирована на функциональные и сценарные свойства: корректность маршрутов API, валидацию версии и режима ответа, правила выбора корпуса, поиск фрагментов, контроль версии, повторное ранжирование, проверку достаточности подтверждений, сохранение истории и работу пользовательского интерфейса.

Программа испытаний построена в соответствии с требованиями, сформулированными в первой главе, и проектными решениями второй главы. Она включает автоматические тесты, сценарные проверки пользовательских функций и сопоставление результатов с требованиями. Количественные метрики качества ранжирования, такие как Recall@K, MRR, nDCG или accuracy, в работе не рассчитывались, поскольку для них требуется отдельный размеченный набор запросов и эталонных фрагментов~\cite{manning2008,beizer1990,myers2011}.

Общая программа испытаний приведена в таблице~\ref{tab:testing-program}. Каждый пункт связан с реализованными файлами практической части и не предполагает ручного изменения состояния корпуса или рабочей базы данных во время проверки отчета.

\begingroup
\centering
\small
\setstretch{1}\selectfont
\setlength{\tabcolsep}{3pt}
\begin{longtable}{|P{0.22\textwidth}|P{0.36\textwidth}|P{0.32\textwidth}|}
\caption{Программа испытаний веб-приложения}\label{tab:testing-program}\\
\hline
\tableheadcell{Направление проверки} & \tableheadcell{Проверяемые свойства} & \tableheadcell{Средство проверки} \\ \hline
\tableheadcell{1} & \tableheadcell{2} & \tableheadcell{3} \\ \hline
\endfirsthead
\multicolumn{3}{l}{\tablepartlabel{окончание таблицы \thetable}}\\
\hline
\tableheadcell{1} & \tableheadcell{2} & \tableheadcell{3} \\ \hline
\endhead
\endfoot
Маршруты API & Доступность диагностики, версий, пользовательского запроса и истории. & \makecell[l]{\texttt{tests/test\_api.py},\\\texttt{tests/test\_history\_api.py}}. \\ \hline
Валидация входных данных & Неподдерживаемая версия, недопустимый режим ответа, короткий вопрос. & Pydantic-схемы и API-тесты. \\ \hline
Выбор корпуса & Официальный корпус для краткого и подробного режимов; оба корпуса для обучающего режима. & \makecell[l]{\texttt{tests/test\_orchestration\_}\\\texttt{rules.py},\\\texttt{tests/test\_reranking.py}}. \\ \hline
Подготовка корпуса & Загрузка локального HTML, обработка markdown-документов, разбиение на фрагменты и проверка покрытия локального официального корпуса. & Тесты загрузчика, покрытия корпуса, парсера и разбиения. \\ \hline
Контроль версии & Исключение официальных источников другой версии и ссылки \texttt{/docs/current/}. & \makecell[l]{\texttt{tests/test\_retrieval\_}\\\texttt{guard.py}}. \\ \hline
Повторное ранжирование & Приоритет официального корпуса и тематических якорей, ограничение дополнительного корпуса. & \texttt{tests/test\_reranking.py}. \\ \hline
Проверка подтверждений & Безопасный отказ при слабом официальном контексте. & \makecell[l]{\texttt{tests/test\_groundedness\_}\\\texttt{guard.py}}. \\ \hline
Генерация и стабилизация результата & Удаление источников из текста ответа, стабилизация обучающей структуры, детерминированные SQL-сценарии. & \makecell[l]{\texttt{tests/test\_generation\_}\\\texttt{safety.py}}. \\ \hline
Пользовательский интерфейс & Три режима ответа, отсутствие переключателя корпуса, история, SQL-блоки. & \makecell[l]{\texttt{tests/test\_frontend\_}\\\texttt{answer\_modes.py}}. \\ \hline
\end{longtable}
\par\endgroup

\section{Критерии проверки}

Критерии проверки определяются не абсолютной точностью ответа, а соответствием поведения системы требованиям. Такой выбор связан с характером ВКР: практическая часть реализует прикладную систему, а не проводит масштабное экспериментальное сравнение поисковых моделей. Критерии перечислены в таблице~\ref{tab:testing-criteria}.

\Needspace{35mm}
\begingroup
\centering
\small
\setstretch{1}\selectfont
\setlength{\tabcolsep}{4pt}
\begin{longtable}{|P{0.30\textwidth}|P{0.58\textwidth}|}
\caption{Критерии проверки результатов}\label{tab:testing-criteria}\\
\hline
\tableheadcell{Критерий} & \tableheadcell{Условие выполнения} \\ \hline
\endfirsthead
\multicolumn{2}{l}{\tablepartlabel{продолжение таблицы \thetable}}\\
\hline
\endhead
\endfoot
Корректность маршрутов API & Реализованные маршруты возвращают структуры данных, соответствующие Pydantic-схемам. \\ \hline
Корректность версии & Неподдерживаемые версии отклоняются, а официальные источники другой версии исключаются из результата. \\ \hline
Корректность режимов & Краткий и подробный режимы используют официальный корпус; обучающий режим допускает дополнительный корпус как вспомогательный. \\ \hline
Трассируемость & Результат содержит список источников или безопасный отказ при недостатке подтверждений. \\ \hline
Безопасный отказ & Система не формирует уверенный ответ при слабом официальном контексте. \\ \hline
История запросов & История содержит вопрос, версию, режим, статус, задержку, результат и источники. \\ \hline
Интерфейс & Главная страница поддерживает ввод вопроса, выбор версии, три режима ответа, источники и недавнюю историю. \\ \hline
\end{longtable}
\par\endgroup

\section{Среда тестирования}

Проверка выполнялась в директории практической части \texttt{/home/arz/Desktop/RAG\_postgres}. Для сохранения режима чтения при запуске тестов использовались переменная \texttt{PYTHONDONTWRITEBYTECODE=1} и отключение поставщика кэша pytest. Это не меняет смысл команды тестирования, но предотвращает создание новых файлов bytecode и кэша из-за самого запуска.

Среда тестирования соответствует зависимостям проекта: Python 3.12, FastAPI 0.116.1, SQLAlchemy 2.0.43, Alembic 1.16.5, pgvector 0.4.1, Pydantic 2.11.7, httpx 0.28.1, Beautiful Soup 4.13.5, sentence-transformers 3.0.1 и pytest 8.4.2. В тестовой конфигурации автоматическая фикстура устанавливает \texttt{APP\_ENV=test}, \texttt{EMBEDDING\_PROVIDER=hashing}, \texttt{EMBEDDING\_DIMENSION=256}, чтобы тесты не зависели от загрузки реальной модели построения векторных представлений.

Контейнерная среда проекта содержит сервисы \texttt{postgres}, \texttt{postgres\_test} и сервис серверной части \texttt{backend}. Здесь \texttt{backend} является техническим именем сервиса в Docker Compose. Основной контейнер базы данных использует образ \texttt{pgvector/pgvector:pg16}. Интеграционные тесты рассчитаны на отдельную тестовую базу данных и пропускаются, если безопасная переменная \texttt{TEST\_DATABASE\_URL} не задана.

\section{Автоматические тесты}

Автоматические тесты размещены в каталоге \texttt{tests}. В проекте выделены модульные, API, сценарные и интеграционные проверки. Основной набор тестов не требует реальной рабочей базы данных и сети: маршруты API проверяются через TestClient и имитированные сервисы, а правила поиска и ранжирования проверяются на синтетических объектах.

Состав автоматических тестов приведен в таблице~\ref{tab:testing-auto-tests}. Она отражает фактический набор файлов тестов в практической части.

\begingroup
\centering
\small
\setstretch{1}\selectfont
\setlength{\tabcolsep}{3pt}
\begin{longtable}{|P{0.28\textwidth}|P{0.38\textwidth}|P{0.24\textwidth}|}
\caption{Автоматические тесты практической части}\label{tab:testing-auto-tests}\\
\hline
\tableheadcell{Файл тестов} & \tableheadcell{Проверяемая логика} & \tableheadcell{Тип проверки} \\ \hline
\tableheadcell{1} & \tableheadcell{2} & \tableheadcell{3} \\ \hline
\endfirsthead
\multicolumn{3}{l}{\tablepartlabel{окончание таблицы \thetable}}\\
\hline
\tableheadcell{1} & \tableheadcell{2} & \tableheadcell{3} \\ \hline
\endhead
\endfoot
\texttt{test\_api.py} & Диагностика, валидация версии и режима, краткий, подробный и обучающий ответы, маршрут версий. & API-тесты с моками. \\ \hline
\texttt{test\_history\_api.py} & Возврат режима, версии и структуры истории запросов. & API-тесты. \\ \hline
\makecell[l]{\texttt{test\_frontend\_}\\\texttt{answer\_modes.py}} & Наличие трех режимов, отсутствие пользовательского переключателя корпуса, история, SQL-блоки. & Проверка HTML/JavaScript. \\ \hline
\makecell[l]{\texttt{test\_retrieval\_}\\\texttt{guard.py}} & Исключение официальных источников другой версии и \texttt{/docs/current/}. & Модульные тесты. \\ \hline
\texttt{test\_reranking.py} & Приоритет официального корпуса, использование дополнительного корпуса в обучающем режиме, тематические бонусы и штрафы. & Модульные тесты. \\ \hline
\makecell[l]{\texttt{test\_groundedness\_}\\\texttt{guard.py}} & Принятие сильного официального контекста и отклонение слабого контекста. & Модульные тесты. \\ \hline
\makecell[l]{\texttt{test\_generation\_}\\\texttt{safety.py}} & Очистка ответа от сгенерированного списка источников, стабилизация обучающего результата, детерминированные SQL-ответы. & Модульные тесты. \\ \hline
\makecell[l]{\texttt{test\_query\_}\\\texttt{processing.py}} & Определение намерения, расширение терминов и признаки логической репликации. & Модульные тесты. \\ \hline
\makecell[l]{\texttt{test\_official\_}\\\texttt{loader.py},\\\texttt{test\_official\_}\\\texttt{corpus\_coverage.py},\\\texttt{test\_parser.py},\\\texttt{test\_chunker\_}\\\texttt{quality.py}} & Загрузка HTML, проверка покрытия официального корпуса, очистка документа и разбиение на фрагменты. & Модульные тесты подготовки корпуса. \\ \hline
\makecell[l]{\texttt{test\_reindex\_}\\\texttt{pipeline.py},\\\texttt{test\_bge\_m3\_}\\\texttt{migration.py},\\\texttt{test\_provider\_}\\\texttt{rules.py}} & Порядок переиндексации, параметры BGE-M3 и запрет хеширующего модуля вне тестов. & Модульные и сценарные тесты. \\ \hline
\texttt{tests/integration/*} & Работа с реальной тестовой PostgreSQL+pgvector базой. & Интеграционные тесты, пропускаются без тестовой БД. \\ \hline
\end{longtable}
\par\endgroup

\section{Сценарные проверки пользовательских функций}

Сценарные проверки дополняют автоматические тесты и показывают, как пользовательские функции связаны между собой. Они описаны по поведению интерфейса и серверной части. Перечень сценариев приведен в таблице~\ref{tab:testing-user-scenarios}.

\begingroup
\centering
\small
\setstretch{1}\selectfont
\setlength{\tabcolsep}{3pt}
\begin{longtable}{|P{0.22\textwidth}|P{0.34\textwidth}|P{0.34\textwidth}|}
\caption{Сценарные проверки пользовательских функций}\label{tab:testing-user-scenarios}\\
\hline
\tableheadcell{Сценарий} & \tableheadcell{Входные данные} & \tableheadcell{Ожидаемый результат} \\ \hline
\tableheadcell{1} & \tableheadcell{2} & \tableheadcell{3} \\ \hline
\endfirsthead
\multicolumn{3}{l}{\tablepartlabel{окончание таблицы \thetable}}\\
\hline
\tableheadcell{1} & \tableheadcell{2} & \tableheadcell{3} \\ \hline
\endhead
\endfoot
Краткий ответ & Вопрос, \texttt{pg\_version=16}, \texttt{answer\_mode=short}. & Возвращается поле \texttt{answer} и официальные источники. \\ \hline
Подробный ответ & Вопрос, \texttt{pg\_version=16}, \texttt{answer\_mode=detailed}. & Возвращается поле \texttt{answer}, источники относятся к официальному корпусу. \\ \hline
Обучающий ответ & Вопрос, \texttt{pg\_version=16}, \texttt{answer\_mode=tutorial}. & Возвращается объект \texttt{tutorial} с четырьмя секциями и источниками. \\ \hline
Недопустимый режим & \texttt{answer\_mode=extended}. & Запрос отклоняется валидацией. \\ \hline
Неподдерживаемая версия & \texttt{pg\_version=19}. & Запрос отклоняется как неподдерживаемый. \\ \hline
Недоступность генерации & Отсутствует \texttt{GROQ\_API\_KEY} или некорректна модель. & Возвращается ошибка сервиса, история фиксирует неуспешный запрос. \\ \hline
Просмотр истории & Запрос \texttt{/api/history?limit=50}. & Возвращаются записи с вопросом, версией, режимом, статусом и источниками. \\ \hline
\end{longtable}
\par\endgroup

После выполнения каждого сценария проверялись не только поля ответа, но и связь результата с выбранным режимом. Для краткого режима достаточно наличия текстового ответа и официальных источников выбранной версии. Для подробного режима дополнительно проверяется полнота объяснения в пределах найденного контекста, но без заявления численной оценки качества. Для обучающего режима проверяется структура результата: краткое объяснение, предварительные условия, шаги и замечания. Если часть обучающих секций пуста, интерфейс скрывает ее, а не показывает пользователю техническую заготовку.

Отдельное внимание уделялось отрицательным сценариям. Неподдерживаемая версия и недопустимый режим ответа должны останавливаться до обработки запроса, поскольку дальнейший поиск в таком случае не имеет корректного корпуса или режима. Недоступность генерации проверяется как сбой внешней зависимости: история фиксирует неуспешный запрос, а пользователь получает сообщение об ошибке. Такой сценарий не подтверждает качество ответа, но подтверждает устойчивость пользовательского потока при отказе внешнего сервиса.

Сценарий просмотра истории связывает пользовательский интерфейс с сохранением результата. Проверяется, что запись содержит вопрос, версию, режим, статус и источники, а сама страница истории не требует повторного выполнения запроса. Пользователь может вернуться к ранее полученному ответу и повторно открыть источники, не обращаясь заново к генеративной модели.

\section{Проверка учёта версии PostgreSQL}

Контроль версии проверяется на двух уровнях. Первый уровень связан с валидацией \texttt{AskRequest.pg\_version}: значение должно быть числовой основной версией и входить в список \texttt{16,17,18}. Второй уровень связан с фильтрацией источников после поиска. Тесты из \texttt{test\_retrieval\_guard.py} проверяют, что официальный источник \texttt{/docs/17/} исключается из результата для версии \texttt{16}, а ссылка \texttt{/docs/current/} не принимается как источник версии \texttt{18}.

Эти тесты подтверждают наличие контроля версии, но не утверждают абсолютную полноту поиска. Система может корректно фильтровать найденные источники, однако качество ответа по конкретному вопросу зависит от полноты локального корпуса и релевантности найденных фрагментов.

\section{Проверка безопасного отказа при недостатке подтверждений}

Проверка достаточности подтверждений тестируется на слабом и сильном контексте. В слабом сценарии вопрос относится к логической репликации, но найденный фрагмент связан с другой темой. Тест подтверждает, что метод \texttt{\_has\_sufficient\_evidence} возвращает \texttt{False}. В сильном сценарии официальный фрагмент содержит понятия \texttt{Logical Replication}, \texttt{publication} и \texttt{subscription}, поэтому проверка возвращает \texttt{True}.

В пользовательском сценарии отрицательный результат проверки приводит не к вызову генеративной модели, а к безопасному отказу. Для краткого и подробного режимов возвращается текст о недостатке подтверждённых данных в официальной документации выбранной версии. Для обучающего режима возвращается структура с пустым списком шагов и поясняющими замечаниями.

\section{Проверка интерфейса и истории запросов}

Пользовательский интерфейс проверяется статическими тестами HTML и JavaScript. Тесты подтверждают, что на главной странице есть только три режима ответа: <<Краткий>>, <<Подробный>> и <<Обучающий>>. Отдельный пользовательский переключатель дополнительного корпуса отсутствует. JavaScript формирует структуру запроса из \texttt{question}, \texttt{pg\_version}, \texttt{answer\_mode}.

Также проверены скрытие пустых секций обучающего ответа, кнопка копирования SQL-блока, блок технических деталей ошибки, ограничение боковой истории четырьмя записями и нейтральный статус завершения на странице истории. API истории проверяет наличие режима и версии в ответе.

\section{Результаты запуска тестов}

Фактический запуск тестов выполнен в практической части командой, приведенной в листинге~\ref{listing:testing-command}.

\begin{listingblock}
\captionof{listing}{Команда запуска автоматических тестов}
\label{listing:testing-command}
\begin{minted}[breaklines=true,breakanywhere=true,fontsize=\small]{bash}
PYTHONDONTWRITEBYTECODE=1 .venv/bin/python -m pytest -q -p no:cacheprovider
\end{minted}
\end{listingblock}

Краткий результат запуска приведен в листинге~\ref{listing:testing-result}.

\begin{listingblock}
\captionof{listing}{Результат запуска автоматических тестов}
\label{listing:testing-result}
\begin{minted}[breaklines=true,breakanywhere=true,fontsize=\small]{text}
91 passed, 3 skipped in 1.00s
\end{minted}
\end{listingblock}

Три теста пропущены, поскольку интеграционные тесты требуют безопасно заданную переменную \texttt{TEST\_DATABASE\_URL} и отдельную тестовую PostgreSQL+pgvector базу. Это поведение описано в \texttt{tests/README.md}: если переменная не задана, интеграционные тесты не выполняются. Такой пропуск не является ошибкой автоматического набора, а защищает рабочую базу данных от случайного использования в тестах. Дополнительное представление команды и вывода тестов приведено в приложении Г в листингах~\ref{listing:appendix-testing-command} и~\ref{listing:appendix-testing-result}.

\section{Сопоставление результатов с требованиями}

Сопоставление требований и проверок приведено в таблице~\ref{tab:testing-req-map}. Она показывает, какие тесты или сценарии подтверждают заявленные функции. Если требование проверяется не одним тестом, а группой проверок, указаны основные файлы и сценарии.

Краткая трассировка пунктов задания и технического задания к реализации приведена в приложении~В, таблица~В.1.

\begingroup
\centering
\small
\setstretch{1}\selectfont
\setlength{\tabcolsep}{2pt}
\begin{longtable}{|P{0.09\textwidth}|P{0.23\textwidth}|P{0.39\textwidth}|P{0.19\textwidth}|}
\caption{Сопоставление требований и тестов}\label{tab:testing-req-map}\\
\hline
\tableheadcell{ID} & \tableheadcell{Требование} & \tableheadcell{Проверяющие тесты и сценарии} & \tableheadcell{Результат} \\ \hline
\tableheadcell{1} & \tableheadcell{2} & \tableheadcell{3} & \tableheadcell{4} \\ \hline
\endfirsthead
\multicolumn{4}{l}{\tablepartlabel{окончание таблицы \thetable}}\\
\hline
\tableheadcell{1} & \tableheadcell{2} & \tableheadcell{3} & \tableheadcell{4} \\ \hline
\endhead
\endfoot
FR-01 & Поддержка версий PostgreSQL & \texttt{test\_api.py}, маршрут \texttt{/api/versions}. & Подтверждено. \\ \hline
FR-02 & Подготовка официального корпуса & \makecell[l]{\texttt{test\_official\_loader.py},\\\texttt{test\_official\_corpus\_coverage.py}.} & Подтверждено. \\ \hline
FR-03 & Разбиение документации на фрагменты & \texttt{test\_parser.py}, \texttt{test\_chunker\_quality.py}. & Подтверждено. \\ \hline
FR-04 & Векторные представления и размерность & \makecell[l]{\texttt{test\_bge\_m3\_migration.py},\\\texttt{test\_provider\_rules.py},\\\texttt{/api/health/embeddings}.} & Подтверждено. \\ \hline
FR-05 & Поиск с учётом версии & \texttt{test\_retrieval\_guard.py}, интеграционный тест при наличии тестовой БД. & Частично подтверждено модульными тестами; интеграционные проверки пропущены без тестовой БД. \\ \hline
FR-06 & Повторное ранжирование & \texttt{test\_reranking.py}. & Подтверждено. \\ \hline
FR-07 & Краткий режим & \texttt{test\_api.py}, \texttt{test\_reranking.py}. & Подтверждено. \\ \hline
FR-08 & Подробный режим & \texttt{test\_api.py}, \texttt{test\_reranking.py}. & Подтверждено. \\ \hline
FR-09 & Обучающий режим & \makecell[l]{\texttt{test\_api.py},\\\texttt{test\_generation\_safety.py},\\\texttt{test\_frontend\_answer\_modes.py}.} & Подтверждено. \\ \hline
FR-10 & Проверка подтверждений & \texttt{test\_groundedness\_guard.py}. & Подтверждено. \\ \hline
FR-11 & История запросов & \texttt{test\_history\_api.py}, интеграционный тест истории при наличии тестовой БД. & Подтверждено на уровне API; интеграционная проверка пропущена без тестовой БД. \\ \hline
FR-12 & Служебная диагностика & \texttt{test\_api.py}. & Подтверждено. \\ \hline
FR-13 & Административная переиндексация & \texttt{test\_reindex\_pipeline.py}, интеграционный тест переиндексации при наличии тестовой БД. & Подтверждено модульным тестом; интеграционная проверка пропущена без тестовой БД. \\ \hline
FR-14 & Пользовательский интерфейс & \texttt{test\_frontend\_answer\_modes.py}. & Подтверждено. \\ \hline
NFR & Нефункциональные требования & Совокупность тестов API, интерфейса, контроля версии, истории, поставщиков векторных представлений и проверки подтверждений. & Подтверждено функционально и сценарно. \\ \hline
\end{longtable}
\par\endgroup

\section{Ограничения разработанного решения}

Разработанное приложение не рассматривается как завершённая промышленная платформа. Основные ограничения связаны с полнотой корпуса, доступностью внешней генеративной модели и методикой оценки. Поддерживаются только версии PostgreSQL, заданные в конфигурации и обеспеченные локальным корпусом: 16, 17 и 18. Добавление новой версии требует наличия корпуса HTML-документации, записи версии и переиндексации.

Качество поиска фрагментов зависит от полноты локальных HTML-страниц, корректности очистки документации, параметров разбиения и качества векторных представлений. В работе не рассчитаны количественные метрики качества ранжирования, потому что не подготовлен эталонный набор запросов и релевантных фрагментов. Поэтому выводы о качестве ограничены функциональными и сценарными проверками.

Генерация итогового текста зависит от доступности Groq API, корректности \texttt{GROQ\_API\_KEY} и выбранной модели \texttt{LLM\_MODEL}. При недоступности генерации система возвращает диагностическую ошибку, но не может сформировать новый текстовый ответ. При этом часть детерминированных SQL-сценариев стабилизирована в коде и проверяется автоматическими тестами.

\section{Выводы}

В четвертой главе определены цель, программа и критерии испытаний, описана среда тестирования, приведены автоматические и сценарные проверки, проверены учёт версии PostgreSQL, безопасный отказ, пользовательский интерфейс и история запросов. Фактический запуск тестов завершился результатом \texttt{91 passed, 3 skipped}. Сопоставление с требованиями показало, что заявленные функции подтверждены автоматическими или сценарными проверками, а ограничения решения связаны с отсутствием количественной оценки ранжирования, зависимостью от полноты корпуса и доступностью генеративной модели.
